\documentclass{PPFIT} % English

\hypersetup{
	pdftitle={LLM-Based Translation Pipeline for Functional Coverage Analysis},
	pdfauthor={Pavel Lukl},
	pdfkeywords={DO-178C, LLM, requirement formalization, agentic pipeline, functional coverage}
}

\PPFinalCopy

\PPYear{2025/2026}

\PaperTitle{LLM-Based Translation Pipeline for Functional Coverage Analysis}

\Authors{Pavel Lukl*}
\affiliation{*%
  \href{mailto:xluklp00@stud.fit.vut.cz}{xluklp00@stud.fit.vut.cz},
  \textit{Faculty of Information Technology, Brno University of Technology}}

\Keywords{DO-178C --- LLM --- requirement formalization --- agentic pipeline --- functional coverage}

%\Supplementary{}

\Abstract{
PLACEHOLDER
% What is the problem?
Safety-critical avionics software under DO-178C requires formal coverage analysis of requirements, but requirements are written in natural language. Manual formalization is slow, error-prone, and does not scale.
%
% How is it solved?
This report investigates the use of large language models (LLMs) in a multi-step agentic pipeline to automatically translate natural-language requirements into a formal predicate-based representation. We survey existing approaches to requirement formalization, evaluate four LLM orchestration frameworks, and prototype a pipeline with ambiguity review, entity extraction, formula generation, and structural validation.
%
% What are the results?
We implemented and compared PydanticAI, Haystack, LangGraph, and CrewAI on a set of real avionics requirements using a locally-run 14B parameter model. CrewAI produced the deepest formulas and fastest execution, while PydanticAI offered the least code complexity and best schema evolution support. The model capability, not the framework, was the primary bottleneck for complex requirements.
%
% So what?
The pipeline architecture is viable for automating requirement formalization. This work establishes the foundation for a production system and identifies concrete next steps: deeper semantic validation, verification via roundtrip checking, tool integration, and evaluation with stronger models.
}

\Teaser{
%	\TeaserImage{placeholder.pdf}
}

\begin{document}
\startdocument

% =============================================================================
\section{Introduction}
\label{sec:intro}
% =============================================================================

Safety-critical avionics software must demonstrate that tests adequately
cover the behaviors specified in requirements. This coverage analysis
requires a formal representation of the requirements, yet in practice,
both requirements and test cases are written in natural language. Without
formalization, inconsistencies between requirements and tests can go
undetected, potentially compromising the safety of the final system.
Translating them into a formal form is therefore essential, but it is
currently a manual process: slow, expensive, and prone to human error.
As the number of requirements grows, manual formalization becomes a
bottleneck in the certification workflow.

The broader project addresses this by building an end-to-end system:
natural language requirements and test cases are translated into a formal
representation, the translations undergo thorough validation, and finally
a gap analysis identifies missing coverage between the formalized
requirements and test cases. This report focuses on the first stage:
automating the translation using large language models (LLMs). Other
parts of the project (the formal representation itself, the semantic
validation of translations, and the gap analysis) are being developed
in parallel by other team members.

Rather than relying on a single LLM prompt, we design an agentic pipeline
of several steps, where each step produces inspectable, validated output,
a property essential in safety contexts where traceability of failures is
required. The pipeline takes a natural language requirement as input and
produces a formal predicate tree conforming to a formalism developed
within the project.

We begin by surveying existing approaches to requirement formalization
and the technologies available for building LLM pipelines in this context
(Section~\ref{sec:sota}). Drawing inspiration from these, we draft our
own approach to the translation problem, describe the pipeline
architecture and the design decisions behind it, and present initial
experiments comparing four orchestration frameworks on real avionics
requirements (Section~\ref{sec:approach}). Based on the findings, we
outline concrete directions for future development as the project
continues (Section~\ref{sec:future}).


% =============================================================================
\section{State of the Art}
\label{sec:sota}
% =============================================================================

Translating natural language requirements into formal representations has
traditionally been a manual process requiring domain expertise. Recent
advances in large language models have opened the possibility of
automating this translation, and a growing body of work explores
different strategies. A recent systematic literature review by Beg
et al.~\cite{beg2025survey} identifies 35~studies on the topic published
between 2022 and mid-2025. We survey the relevant approaches below,
organized into four areas: requirement quality and pre-processing
(Section~\ref{sec:sota-quality}), formalization approaches
(Section~\ref{sec:sota-formalization}), output representation and
correctness (Section~\ref{sec:sota-output}), and implementation
technologies (Section~\ref{sec:sota-tech}).

% -----------------------------------------------
\subsection{Requirement Quality and Pre-processing}
\label{sec:sota-quality}

Before a requirement can be formalized, it must be sufficiently
unambiguous and complete. Template-based tools such as
FRET~\cite{giannakopoulou2020fret} achieve precise formalization by
having users write requirements in a restricted natural language with
unambiguous semantics, but require every requirement to be manually
rewritten to conform to the template grammar. NLP-based defect detection
offers an alternative: pattern-based approaches have been applied to
identify ambiguities and inconsistencies in requirements in safety-critical
domains such as railway~\cite{ferrari2018detecting}. More recently, Mahbub
et al.~\cite{mahbub2024gpt4ambiguity} evaluated GPT-4 for automated
detection of ambiguities, inconsistencies, and incompleteness in
industrial requirements, finding that while LLMs show promise for early
defect screening (precision 0.61 for incompleteness), they are less
reliable for ambiguity detection (precision 0.39) and cannot replace
human analysts. These findings suggest that LLM-based ambiguity review is
useful as a pre-processing step that flags potential issues for human
review, but should not autonomously rewrite requirements.

% -----------------------------------------------
\subsection{Formalization Approaches}
\label{sec:sota-formalization}

\paragraph{Interactive decomposition and dynamic prompting.}
Cosler et al.~\cite{cosler2023nl2spec} propose decomposing the
formalization into sub-translations: the LLM maps fragments of the input
to sub-formulas, which serve as a structured intermediate representation
that users can inspect, edit, and correct before the final formula is
assembled. This decomposition makes the model's reasoning transparent at
each step, and the interactive loop lets users resolve ambiguities that
single-pass approaches cannot reliably detect. The framework, called
\textit{nl2spec}, targets LTL and similar temporal logics. Xu
et al.~\cite{xu2024nl2ltl} take the prompting strategy further with
dynamic few-shot prompting (where a small number of solved examples are
included in the prompt to guide the model): instead of using a fixed set
of examples, the system retrieves the most relevant ones from a prompt
library via embedding similarity for each input, and evolves the library
over time by incorporating corrected translations. This achieves up to
94.4\% accuracy without human interaction. With interactive prompt
evolution, accuracy on an out-of-domain dataset improves from 27\% to
78\%. Both approaches demonstrate that careful prompt selection and
human-in-the-loop refinement can compensate for a model's limited domain
knowledge, and that inspectable intermediate outputs are valuable for
catching errors early.

\paragraph{Structured decomposition and agentic pipelines.}
Khot et al.~\cite{khot2022decomposed} show that decomposing complex tasks
into sub-tasks delegated to specialized prompts outperforms single-pass
chain-of-thought, since each prompt can be independently optimized for its
sub-task. Several recent works apply this principle to requirement
formalization, separating semantic reasoning from formal output generation
using structured intermediate representations.
Ma et al.~\cite{ma2025req2ltl} propose REQ2LTL, where the LLM performs
semantic reasoning about the requirement and produces output in a
hierarchical intermediate representation called OnionL, a tree of scopes,
relations, and atomic propositions. A deterministic rule-based engine then
validates this tree and translates it to LTL. The LLM is called with a
single prompt containing the grammar of the intermediate language, a
step-by-step decomposition algorithm as chain-of-thought instructions,
several ground-truth examples of requirement-to-tree translations, and the
requirement itself. This achieves 100\% syntactic correctness and 88.4\%
semantic accuracy on 112 real aerospace requirements, significantly
outperforming other approaches. The framework also supports optional human
validation via automatically generated visual tree diagrams. An ablation
study confirms that removing either the intermediate representation or
the chain-of-thought decomposition degrades accuracy by over 20
percentage points. Yan et al.~\cite{yan2025assertllm} apply a similar
separation in hardware verification: the first LLM converts unstructured
text into structured descriptions using a unified template, the second
generates assertions from these descriptions, and the third validates the
results. Nouri et al.~\cite{nouri2024safety} use the same principle in
the automotive safety domain, breaking the analysis into smaller tasks
with a dedicated prompt for each, where each step's output feeds the next
without modification. Their key finding is that this decomposition
significantly outperforms a single monolithic prompt. Together, these
works establish that separating semantic extraction from formal
translation, with validated intermediate representations between stages,
is an effective strategy for complex formalization tasks.

\paragraph{Fine-tuning.}
Fine-tuning language models on datasets of natural language sentences
paired with their formal translations~\cite{hahn2022formal} can achieve
competitive accuracy, but requires a large training dataset that does not
exist for specialized domains such as avionics.

% -----------------------------------------------
\subsection{Output Representation and Correctness}
\label{sec:sota-output}

A recurring challenge across all formalization approaches is ensuring
that the generated output is not only structurally valid but also
semantically correct. This subsection surveys strategies for improving
output correctness, from constraining the generation process to choosing
the right output format.

\paragraph{Constrained decoding.}
Grammar-Forced Translation (GraFT) by English
et al.~\cite{english2025graft} constrains the language model's token
generation at each step so that only tokens permitted by the target
grammar can be produced. This provides a formal guarantee that every
generated formula is syntactically valid, eliminating the need for
structural retries entirely. The technique involves two steps
transferable to other formalisms: first, a masked language model lifts
atomic propositions from the natural language input (replacing domain
terms with integer placeholders), then a fine-tuned sequence-to-sequence
model generates the formula under grammar constraints. In their
evaluation on temporal logic benchmarks, this improves end-to-end
accuracy by 5.49\% and out-of-distribution accuracy by 14.06\% over
prior methods.

\paragraph{Code as output representation.}
The idea of using code rather than custom notation as the output format
has broader support. Gao et al.~\cite{gao2022pal} show that having LLMs
generate Python programs instead of reasoning directly, then executing
the code, outperforms chain-of-thought on mathematical and symbolic
reasoning tasks by offloading the solution step to an interpreter.
Danso et al.~\cite{danso2026syntax} provide evidence specific to
formalization: in a systematic evaluation of several LLMs on NL-to-LTL
translation, they compare three prompt formats with increasing structural
constraint (minimal, detailed grammar, and a Python interface where each
temporal operator is a Python dataclass constructor). Reformulating the
task as Python code generation improves semantic equivalence from
approximately 24\% to over 70\%, and raises syntactic correctness to
near 100\%. The improvement is attributed to two factors: LLMs have far
more Python training data than LTL notation, and the dataclass
constructors physically prevent certain classes of malformed output.

\paragraph{The syntax--semantics gap.}
Critically, Danso et al.\ show that the dominant source of error is not
syntax but semantics. While models produce syntactically valid formulas
most of the time, they frequently get the \textit{meaning} wrong. The
most common errors are: confusion between the ``always'' (G) and
``eventually'' (F) operators, incorrect nesting of temporal scopes
(placing an operator at the wrong level of the formula tree), and over-
or under-specification where the generated formula is strictly stronger
or weaker than intended. Atomic proposition extraction is relatively
reliable (57.2\% overlap), but operator agreement drops to 36.7\% and
exact structural alignment to 18.7\%. These semantic errors are invisible
to structural validation and can only be detected through formal
equivalence checking, which the authors perform using the NuSMV model
checker as an external oracle. This finding is echoed by
REQ2LTL~\cite{ma2025req2ltl}, whose ablation study shows that structured
decomposition addresses precisely these semantic errors, improving
accuracy by over 20 percentage points. LLM compliance with complex JSON
schemas also degrades sharply as schema depth
increases~\cite{geng2025jsonschemabench}, motivating either simpler
per-step schemas or alternative output representations.

% -----------------------------------------------
\subsection{Technologies}
\label{sec:sota-tech}

Building a formalization pipeline requires concrete choices about how the
steps are orchestrated and how the LLM is accessed. Several open-source
frameworks have emerged for this purpose, and a recent comparison of ten
such frameworks on the same task found significant differences in
implementation complexity and
reliability~\cite{ahmed2026ddl2propbank}. PydanticAI provides the most
direct integration with Pydantic output schemas, with automatic retry on
validation failure built in. Haystack models pipelines as directed acyclic
graphs of typed components, with native support for serialization to YAML.
LangGraph represents pipelines as state graphs with conditional edges and
supports cycles, which naturally accommodates loops that retry on
validation failure. CrewAI takes a role-based approach, defining agents
with goals and backstories that collaborate on tasks. Other notable
frameworks include DSPy, which optimizes prompts programmatically, and
Instructor, which wraps function-calling APIs with Pydantic validation.
The choice of framework determines how easily the pipeline can be
inspected at intermediate steps, how schema changes propagate, and how
validation errors are surfaced, all of which directly affect auditability.
The LLM itself can be accessed either through cloud APIs (such as Claude
or GPT-4) or by running open-weight models locally using tools like
Ollama. Most orchestration frameworks support both options, allowing the
same pipeline to switch between them with minimal code changes.

% -----------------------------------------------
\subsection*{Summary}

The surveyed literature shows that LLMs can translate natural language
requirements into formal representations with increasing accuracy, and
that the most effective approaches share several methodological patterns:
task decomposition into smaller, individually verifiable steps; structured
intermediate representations between stages; feedback-driven refinement
through human interaction, dynamic prompt evolution, or automated retry;
and careful choice of output representation, with Python code
outperforming custom notation on frontier
models~\cite{danso2026syntax,gao2022pal}. A recurring finding is that
syntactic correctness is largely solved, while the real challenge lies in
semantic correctness, particularly mis-scoped temporal operators and
incorrect nesting~\cite{danso2026syntax,ma2025req2ltl}. The closest
existing work is REQ2LTL~\cite{ma2025req2ltl}, which targets aerospace
requirements and uses a hierarchical intermediate representation
conceptually similar to our formalism. However, their output is standard
LTL and their translation step is a fixed rule-based engine, whereas our
formalism is domain-specific and still evolving, requiring the
translation step itself to use LLM generation with schema validation. Our
work builds on the same principles demonstrated by REQ2LTL and the other
surveyed approaches, applied to this different setting.


% =============================================================================
\section{Our Approach}
\label{sec:approach}
% =============================================================================

% -----------------------------------------------
\subsection{Research Questions}
\label{sec:rqs}

Based on the patterns and open questions identified in the literature, we
investigate four research questions:

\begin{itemize}
\item \textbf{RQ1: How should the formalization task be decomposed?}
The literature shows that both single-call chain-of-thought
decomposition~\cite{ma2025req2ltl} and multi-call
pipelines~\cite{nouri2024safety,khot2022decomposed} outperform monolithic
prompts, but does not establish which is preferable for a given formalism
and model capability level.

\item \textbf{RQ2: How does output representation affect formalization
accuracy?} Danso et al.~\cite{danso2026syntax} and Gao
et al.~\cite{gao2022pal} report that Python code output significantly
outperforms custom notation on frontier models. We test whether this
finding holds for smaller local models and for our specific formalism.

\item \textbf{RQ3: How does model capability affect architecture
choices?} Cloud frontier models and local open-weight models differ
substantially in reasoning ability. We investigate whether the optimal
pipeline architecture depends on the model used.

\item \textbf{RQ4: Which orchestration framework best fits this use
case?} Ahmed et al.~\cite{ahmed2026ddl2propbank} show significant
differences in framework complexity and reliability. We compare four
frameworks on the same formalization task.
\end{itemize}

% -----------------------------------------------
\subsection{The Formalism}
\label{sec:formalism}

The target representation is a domain-specific predicate formalism
designed for expressing avionics software requirements and test cases in
a form suitable for mechanical coverage analysis. We summarize the
aspects relevant to the pipeline.

A formalized requirement is a triple $R$ consisting of a
\textit{flavour} (Discrete or Continuous, fixing the time model), a set
of \textit{entities}, and a \textit{constraint} (a temporal construct
expressing the required behavior). Entity types include Signal, Storage, Event, EventTrigger,
Channel, State, Value, Abstract, and Set, each with a defined set of
permitted operations (e.g., \texttt{fire} for EventTrigger,
\texttt{write} for Storage). Temporal constructs such as Always,
Eventually, Causes, CausesWithin, Sequence, and others quantify over
time, and predicates within them reference entity values, event
occurrences, and interval-based computations.

The formalism can be mapped to LTL and similar temporal logics, which
makes the LTL-focused literature from Section~\ref{sec:sota} directly
informative. Our pipeline targets the formalism itself as its output,
analogous to how REQ2LTL~\cite{ma2025req2ltl} targets OnionL as its
intermediate representation. Once the formalism stabilizes and a mapping
to temporal logics is defined, a deterministic rule-based translation to
LTL could be added as a downstream step, mirroring REQ2LTL's
architecture.

% -----------------------------------------------
\subsection{Pipeline Architecture}
\label{sec:architecture}

Following the decomposition principle established in the
literature~\cite{khot2022decomposed,nouri2024safety}, our pipeline
breaks the formalization into four steps, each producing inspectable,
validated output:

\begin{enumerate}
\item \textbf{Ambiguity Review.} The LLM flags potential ambiguities in
the requirement text against a checklist (temporal ambiguity, scope of
negation, missing causes, vague language). Following Mahbub
et al.'s~\cite{mahbub2024gpt4ambiguity} finding that LLMs are unreliable
for autonomous ambiguity resolution, this step flags issues for human
review but does not rewrite the requirement. The ambiguity notes are
passed as context to all subsequent steps.

\item \textbf{Flavour Extraction.} The LLM determines whether the
requirement uses discrete time (step-based behavior, register operations)
or continuous time (real-time durations, deadlines). This is a simple
classification that sets the time model for the rest of the pipeline.

\item \textbf{Entity Extraction.} The LLM identifies all entities
referenced in the requirement, producing a structured intermediate
representation: each entity is a tuple of (name, type, modifiers,
description). This step serves the same architectural role as the
intermediate representation in
REQ2LTL~\cite{ma2025req2ltl} and
AssertLLM~\cite{yan2025assertllm}: it separates semantic extraction
from formal output generation and provides an inspection point where
entity types and modifiers can be validated before the constraint is
built.

\item \textbf{Constraint Formalization.} Given the entities, flavour,
and ambiguity notes, the LLM constructs the temporal constraint. The
prompt includes a step-by-step decomposition algorithm (inspired by
REQ2LTL's BUILDING-ONIONL~\cite{ma2025req2ltl}): identify the top-level
temporal scope, decompose the condition into sub-predicates, decompose
the effect, and normalize entity references. If structural validation
fails, the error is fed back and the step retries automatically.
\end{enumerate}

Structural validation checks that every entity referenced in the
constraint was declared in step~3, that operations are compatible with
entity types, and that the output conforms to the Pydantic schema. This
catches structural errors but not semantic ones: as Danso
et al.~\cite{danso2026syntax} show, a formula can be structurally valid
while meaning something entirely different from the requirement. Semantic
validation is a separate concern outside the scope of this work.

% -----------------------------------------------
\subsection{Preliminary Experiments}
\label{sec:experiments}

We conducted four sets of preliminary experiments to inform the research
questions. Due to the exploratory nature of this phase, sample sizes are
small and results should be interpreted as directional findings, not
definitive answers. Each experiment's setup (model, schema version,
requirements tested) is noted explicitly, as these varied across
experiments.

\paragraph{RQ4: Framework comparison.}
We implemented the pipeline in four orchestration frameworks (PydanticAI,
Haystack, LangGraph, CrewAI) and ran 48~tests (4~frameworks $\times$
4~requirements $\times$ 3~runs) using a locally served qwen2.5:14b model
and an earlier, simpler version of the formalism schema. Three frameworks
achieved 92\% structural success rate; LangGraph reached 75\%. CrewAI was
consistently fastest (25.5s average vs 50--66s for others) and produced
the deepest formula trees. Investigation via HTTP request interception
revealed that CrewAI omits the \texttt{response\_format: json\_object}
parameter that the other frameworks send, but reproducing this change in
PydanticAI did not reliably replicate the advantage. PydanticAI required
the least code (${\sim}$140 lines vs ${\sim}$200--220) and offered the
best schema propagation: changing the Pydantic model automatically
updates the LLM's schema context.
\paragraph{RQ2: Output representation.}
To test Danso et al.'s~\cite{danso2026syntax} finding that Python output
improves formalization accuracy, we ran 10~requirements $\times$ 3~runs
in both JSON and Python output modes using phi4:14b (a local model) and
the updated formalism schema. The LLM received the same formalism context
in both cases; only the output format instructions differed. JSON
achieved 90\% structural success (27/30) while Python achieved 70\%
(21/30). Python errors included inventing function names, wrong argument
counts, and confusing API levels. JSON was also faster (17.3s vs 20.2s
average). This contradicts Danso et al.'s finding, but their evaluation
used frontier models (GPT-4o, Claude, Gemini). Our result suggests that
the Python advantage is capability-dependent: it emerges when the model
is fluent enough in Python to benefit from the structural scaffolding,
but a smaller model makes more errors with an unfamiliar API than with
JSON discriminators it has seen in training.

\paragraph{RQ1: Task decomposition.}
We compared two architectures: the multi-call pipeline (separate LLM
calls for flavour, entities, and constraint with validation between
steps) and a single-call chain-of-thought approach (one prompt containing
a step-by-step decomposition algorithm, inspired by
REQ2LTL~\cite{ma2025req2ltl}). Both used identical decomposition logic;
only the execution model differed. On phi4:14b, the pipeline achieved 75\%
success vs 50\% for CoT (4~requirements, 1~run). On Gemini~3.5 Flash (a
cloud model), both achieved 100\% (4~requirements, 1~run). Ground truth
comparison revealed that CoT produced better temporal precision on complex
requirements (correctly using \texttt{LastOcc} to reference values at
past event times), while the pipeline produced more thorough entity
extraction. An informal test with a frontier model (Claude) on the
hardest requirement produced the closest match to ground truth of any
configuration, correctly capturing temporal lookups that all other setups
missed. These results suggest that weaker models benefit from the
pipeline's simpler per-call schemas, while stronger models benefit from
CoT's full-context reasoning.

\paragraph{RQ3: Model capability.}
Across all experiments, model capability was the strongest predictor of
output quality. The local phi4:14b model achieved 75--90\% structural
success on simple to medium requirements but failed consistently on the
most complex requirement (Req~9: multi-condition timed toggle with event
counting and interval construction). Gemini~2.5 Flash and 3.5~Flash
succeeded on this requirement in both pipeline and CoT modes. The
semantic quality gap was even larger: cloud models correctly used
advanced constructs such as event occurrence counting with open intervals
and value lookups at past event times, which local models either
simplified away or produced incorrectly. Reasoning-specific model
variants were incompatible with framework-managed structured output due
to embedded thinking tags that interfere with JSON parsing.

% -----------------------------------------------
\subsection{Discussion}
\label{sec:discussion}

The preliminary experiments confirm several findings from the literature
and reveal one notable divergence. Task decomposition improves
formalization quality regardless of whether it is implemented as a
multi-call pipeline or as chain-of-thought instructions within a single
call, confirming the core finding of Khot
et al.~\cite{khot2022decomposed} and
REQ2LTL~\cite{ma2025req2ltl}. The choice between the two
architectures depends on model capability: smaller models need the
simpler per-step schemas that a pipeline provides, while larger models
benefit from seeing the full requirement context in a single call.

The finding that JSON outperforms Python on a 14B local model (RQ2) is
the clearest divergence from the literature. Danso
et al.~\cite{danso2026syntax} measured semantic equivalence improvements
of nearly threefold with Python output, but their evaluation used
frontier models. Our result suggests that the advantage is not inherent
to the output format but depends on the model's fluency with the
representation: a model that has seen far more JSON schemas in training
than custom Python APIs will perform better with JSON. This has practical
implications for deployment in environments where cloud models are not
available.

All findings reported above are based on small sample sizes (1--3 runs
per configuration, 4--10 requirements) and should be treated as
preliminary observations rather than confirmed results. Confirming these
trends requires more comprehensive testing: larger requirement sets,
multiple runs per configuration for statistical significance, a wider
range of models, and systematic ground-truth evaluation. The resources
required for such testing, particularly access to frontier cloud models
with sufficient API quotas, were not available during this phase of the
project.

Several design questions also remain open. The degree of human involvement
(fully automated vs human-in-the-loop review, as in
nl2spec~\cite{cosler2023nl2spec} and REQ2LTL~\cite{ma2025req2ltl})
depends on the deployment context and the level of trust required in the
output. Dynamic few-shot example selection (Xu
et al.~\cite{xu2024nl2ltl}) could improve accuracy by retrieving the
most relevant examples for each input but has not yet been tested.
Semantic validation beyond structural checks remains the fundamental open
problem: as Danso et al.~\cite{danso2026syntax} show, a structurally
valid formula can be semantically incorrect, and detecting such errors
requires formal equivalence checking or human review.


% =============================================================================
\section{Future Work}
\label{sec:future}
% =============================================================================
% Critical — the project continues. Write so that future-me can pick up where I left off.
%
% TODO: Cover all of the following as a continuous section (no subsections).
%
% Pipeline improvements:
% - Deeper semantic validation: operator validity, time variable consistency,
%   argument count/type checking, modifiers validation
% - Entity extraction: unique naming enforcement, typing consistency,
%   validate before passing to formula step
% - Ambiguity review: only flag, never rewrite (rewriting introduced errors)
% - Schema gaps: Constant.value is float-only, no arithmetic expressions
%
% Verification and batch processing:
% - Roundtrip verification via cosine similarity (from pipeline diagram)
% - Operational guardrails generation
% - Entity deduplication across requirements
% - Gap analysis: formalized requirements vs formalized test cases
% - Coverage metrics
%
% Tool integration:
% - Syntax checker, domain glossary, requirements database
% - PydanticAI has strongest tool support
%
% Model and deployment:
% - Cloud model testing (Claude, GPT-4) to establish quality ceiling
% - Local vs cloud tradeoff for safety-critical environments
% - Fine-tuning on domain-specific data
% - Production: latency, cost, auditability
% - Schema evolution as formalism matures
% - Extended evaluation: all 10 requirements, human evaluation,
%   comparison with manual formalization
%
% Alternative output representations:
% - Investigate translating to C++ code instead of/alongside the predicate formalism
%   — LLMs are heavily trained on code and may produce more reliable output
%   in a programming language than in a custom formal notation
% - Could serve as an intermediate representation or as a verification target


% =============================================================================
\section{Conclusion}
\label{sec:conclusion}
% =============================================================================

% [Summary]
% TODO: What we set out to do — investigate LLM pipelines for NLR formalization.
% What we did — surveyed state of the art, designed a pipeline, prototyped in
% four frameworks, ran 48 comparison tests, investigated model differences.

% [Key findings]
% TODO:
% - The pipeline architecture works for simple-to-medium requirements
% - Model capability is the bottleneck, not the framework
% - CrewAI and PydanticAI continue as candidates
% - Local 14B model is viable for development, insufficient for complex requirements
% - Prompt engineering significantly affects output quality

% [What this enables]
% TODO: This work establishes the foundation — pipeline architecture, schemas,
% prompts, and empirical data — for building a production requirement
% formalization system in a DO-178C context.

% [Next semester]
% TODO: Priorities:
% - Deeper validation (operators, time variables, arguments)
% - Roundtrip verification
% - Tool integration (syntax checker, glossary, requirements DB)
% - Cloud model evaluation
% - Scale to full requirement set with human evaluation


\phantomsection
\ifenglish
    \bibliographystyle{bib-styles/enplain}
\fi

\bibliography{2019-PPFIT-LLMPipeline-bib}

\end{document}
